% Independent continuation by /root/frobenius, after C.4m--n.
% The final scope paragraph supersedes the earlier intermediate scope paragraphs.
\subsection{The Prym--Tyurin improvement for arbitrary threefolds}
\label{sec:c4-pt-six}

\paragraph{C.4o (Proved: six times the identity on every principally polarized sixfold).}
Let $(B,\Xi)$ be a principally polarized abelian variety of dimension
six over an algebraically closed field of arbitrary characteristic.
There is an integral curve cycle $\alpha$ such that
\begin{equation}\label{eq:c4-six-minimal}
       \operatorname{cl}(\alpha)=6\,\frac{\xi^5}{5!},
       \qquad \xi=c_1(\Xi),
\end{equation}
and consequently $6\in I_J(B)$.

\emph{Proof.}
Over $\mathbb C$, the integral-cycle assertion is exactly
\citet[Corollary 2, p.~165]{frC4ADFIO2020}.
Their Theorem 1 on the same page proves dominance of the
Prym--Tyurin map of exponent six. The meaning of that exponent is
the identity $j^*\Theta_C\equiv6\Xi$ for the map
$j:B\hookrightarrow J(C)$ on a general member, as specified on
p.~164. Polarized adjunction gives $j^\dagger j=[6]$;
the corollary includes the specialization to \emph{every} complex
principally polarized sixfold. In particular, we do not replace the
word ``every'' by an unsupported assertion about the generic member.

We justify transfer from $\mathbb C$ to an arbitrary algebraically
closed field $F$ of characteristic zero in a form that also preserves
the integral coefficient. Define $(B,\Xi)$ over a finitely generated
subfield $K\subset F$, and choose an embedding $K\hookrightarrow
\mathbb C$. Write the curve furnished by the complex theorem as
\[
             \alpha=\sum_a m_a(f_a)_*[C_a],
\]
with smooth connected projective $C_a$, chosen base points, and
induced homomorphisms $\bar f_a:J(C_a)\to B_{\mathbb C}$.
Its class implies the identity
\begin{equation}\label{eq:c4-six-spread-identity}
 \sum_a m_a\bar f_a\lambda_a^{-1}
                  \widehat{\bar f_a}\lambda_B=[6]_B
\end{equation}
by the Matsusaka calculation in C.4i. The finitely many curves,
base points, maps, and identity spread to a nonempty scheme $S$ of
finite type over $K$, after shrinking so that the curves are smooth
and projective. Their relative Jacobians and principal polarizations
spread as well. A closed point of $S$ has residue field finite over
$K$, which embeds over $K$ into $F$. Its fibre therefore supplies
the same data over $F$, still satisfying
\eqref{eq:c4-six-spread-identity}.

As in the alternating-pairing argument in C.4m, the associated
curve on $B_F$ has class $6\xi^5/5!$: the symmetric map
$\widehat B\to B$ defined by the left side before composition
with $\lambda_B$ is $6\lambda_B^{-1}$. Thus this argument transfers
the full curve-class assertion, not merely an unpolarized
factorization whose adjoint relation might have been lost.

If $\operatorname{char}F=p>0$, lift $(B,\Xi)$ over $W(F)$ using
\citet[Corollary 2.4.2 and its proof, p.~287]{frC4Oort1971}.
The geometric generic fibre has characteristic zero, so the preceding
paragraph supplies its integral curve class. We descend the finite
cycle data to a finite extension of the fraction field, close the
cycle, and specialize it as in C.4m. The compatibility in
\citet[Example 20.3.5, p.~400]{frC4Fulton1998} yields
\eqref{eq:c4-six-minimal} in rational $\ell$-adic cohomology for
$\ell\ne p$, hence numerically. The integral Matsusaka formula,
with $\chi(\Xi)=1$, gives $6\in I_J(B)$.
This argument includes $p=2$ and $p=3$: the cycle remains integral
even when the coefficient six is divisible by the characteristic.
\hfill$\square$

\paragraph{C.4p (Proved: the threefold obstruction divides six).}
For every abelian threefold $A$ over an algebraically closed field,
\begin{equation}\label{eq:c4-threefold-six}
 \boxed{\qquad s_J(A)\mid\gcd(6,\delta(A)).\qquad}
\end{equation}
In particular, Question 2 holds for $A$ if, for each of
$\ell=2,3$, there is an ample line bundle $L_\ell$ with
$\ell\nmid\chi(L_\ell)$. The two line bundles may be different.
For a complex Picard-rank-one threefold with primitive polarization
of type $(1,d_2,d_3)$, the sufficient condition is
$\gcd(d_2d_3,6)=1$.

\emph{Proof.}
For a sixfold $X$, principalizing an ample line bundle $L$ and
applying C.4o gives
\[
                  6\chi(L)\in I_J(X).
\]
The same ample-gcd argument as before gives
\[
                         s_J(X)\mid6\delta(X).
\]
For $X=A\times\widehat A$, the Poincar\'e line bundle has Euler
characteristic $-1$, so $\delta(X)=1$. Restriction and projection
give $6\in I_J(A)$. C.4j gives
$\delta(A)\in I_J(A)$; taking the gcd proves
\eqref{eq:c4-threefold-six}. The local condition at two and three
means exactly that this gcd is one. Combining the finite
factorizations by B\'ezout and using C.4h produces an actual direct
summand. For Picard rank one, the Riemann--Roch calculation gives
$\delta(A)=\chi(L)=d_2d_3$.\hfill$\square$

The gcd construction is finite here too. Choose an ample $H$ on
$X=A\times\widehat A$ and $T$ so large that
$\mathcal P_A\otimes H^{T+j}$ is ample for $0\le j\le6$.
For $p(t)=\chi(\mathcal P_A\otimes H^t)$, the identity
\[
 -1=p(0)=\sum_{j=0}^{6}(-1)^j
          \binom{T+j-1}{j}\binom{T+6}{6-j}\,p(T+j)
\]
combines seven available factorizations of $6p(T+j)$ into one
of $-6$, and changing the sign of the return homomorphism gives
one of $6$.

\paragraph{The remaining low-dimensional cases.}
The strongest bounds established in this contribution are now
\[
\begin{array}{c|c|c}
 \dim A & \text{proved bound for }s_J(A)
        &\text{possible residual primes here}\\\hline
 2 & \gcd(2,\delta(A)) & 2\\
 3 & \gcd(6,\delta(A)) & 2,3.
\end{array}
\]
Thus the earlier possible prime three for surfaces and prime five
for threefolds have been removed. We have neither proved the
unrestricted direct-summand assertion for all remaining varieties
nor produced a counterexample among them. For example, a
Picard-rank-one surface of type $(1,2)$ remains untreated here;
type $(1,3)$ is affirmative. For threefolds, type $(1,1,5)$ is
affirmative, while the present bound alone does not decide
type $(1,1,3)$ or type $(1,2,2)$.

The bound for principally polarized sixfolds is sharp on a very
general complex member: \citet[Theorem 1.2]{frC4EGFSOptimality2025}
prove that the smallest algebraic multiple of its minimal curve
class is six. This does not give a counterexample for a threefold,
because the special sixfold $A\times\widehat A$ is not a very
general principally polarized sixfold. Likewise, the evenness
theorem for very general principally polarized fourfolds does not
decide the remaining nonprincipal surface cases. These distinctions
are essential to the exact integral question.
